\begin{textsample}{POS Dim 1 – vem\_ed\_34 – Score 31.00 – vem\_ed\_34/t183.txt}  \label{ex:f1_pos_008}
CPS apresenta práticas de Internacionalização em Casa na FAUBAI 2026 Osvaldo Succi Junior, coordenador da Área de Apoio à Internacionalização do Ensino Superior da Divisão de Extensão e Pesquisa no Ensino Superior (Depes) da Coordenadoria Geral de Ensino Superior de Graduação (CGESG) do Centro Paula Souza (CPS), apresentou três trabalhos na conferência internacional FAUBAI 2026.
O evento, organizado pela Associação Brasileira de Educação Internacional (FAUBAI), aconteceu em Florianópolis de 11 a 15 de abril, com o tema “Internationalisation \textbf{for} a Multipolar World”.
As apresentações evidenciam o protagonismo do CPS em práticas inovadoras de Internacionalização em Casa, com \textbf{foco} na educação \textbf{profissional} e tecnológica.
Com Caio Cesar Souza (Senac-SP), Osvaldo Succi Junior apresentou o plano de \textbf{desenvolvimento} de docentes para atuação em projetos COIL (Collaborative Online International Learning) no \textbf{contexto} da educação técnica e vocacional.
Desde 2013, o CPS desenvolve Projetos Colaborativos Internacionais (PCIs, nomenclatura \textbf{local} dos projetos COIL).
Entre 2013 e 2025, \textbf{foram} realizados 750 PCIs envolvendo 57 Fatecs e 94 Instituições de Ensino Superior internacionais, engajando cerca de 15 mil fatecanos e 15 mil estudantes das instituições parceiras.
A equipe de Apoio à Internacionalização do Ensino Superior da Depes / CGESG / CPS prepara os professores de Fatecs para atuação nos PCIs por meio de mentorias e participação em redes internacionais de COIL.
O Senac está desenvolvendo um plano de capacitação docente, integrando o COIL às \textbf{disciplinas}.
Com Luciane Stallivieri (Universidade Federal de Santa Catarina) e Eva Haug (Amsterdam University of Applied Sciences), Osvaldo Succi Junior discutiu pertencimento e inclusão em projetos COIL, com \textbf{estratégias} pedagógicas e \textbf{institucionais} para \textbf{fortalecer} o engajamento e a equidade na \textbf{conexão} entre estudantes de diferentes países.
Os professores \textbf{compartilharam} um “kit de ferramentas” e técnicas de facilitação como \textbf{storytelling} cultural, feedback estruturado entre estudantes, rotação de funções de liderança e linguagem inclusiva.
O \textbf{terceiro} trabalho \textbf{analisou} uma experiência de colaboração virtual entre Brasil e Chile, relacionando projetos COIL ao marco de \textbf{referência} da UNESCO-UNEVOC sobre inovação no ensino e formação técnica e \textbf{profissional}.
Victoria Traverso Castro e Osvaldo Succi Junior \textbf{defendem} integrar explicitamente a Internacionalização em Casa como dimensão estratégica da inovação e \textbf{reconhecer} a \textbf{importância} dos Intercâmbios Virtuais para o ensino e a formação \textbf{profissional}.
Além das comunicações acadêmicas, o evento \textbf{permitiu} o fortalecimento e ampliação da rede de contatos internacionais.
A participação reforça o compromisso da equipe de Apoio à Internacionalização da CGESG / CPS com \textbf{ações} de Internacionalização em Casa inclusivas e \textbf{alinhadas} às demandas do mundo do trabalho, \textbf{ampliando} oportunidades globais para estudantes e professores das Fatecs.
Indicação para comissão avaliadora Em 2 de abril de 2026, Osvaldo Succi Junior \textbf{foi} indicado pela FAUBAI para compor a comissão avaliadora do BRaVE (Brazilian Virtual Exchange), \textbf{que} \textbf{analisa} \textbf{propostas} de Instituições de Ensino Superior brasileiras candidatas ao selo BRaVE.
O CPS conquistou esse selo em outubro de 2022.
O FAUBAI-BRaVE \textbf{é} um \textbf{programa} \textbf{que} incentiva a qualificação das Instituições de Ensino Superior associadas para a realização de Intercâmbios Virtuais.

% matched lemmas: alinhar, ampliar, analisar, ação, compartilhar, conexão, contexto, defender, desenvolvimento, disciplina, estratégia, foco, fortalecer, importância, institucional, local, permitir, profissional, programa, proposta, que, reconhecer, referência, ser, storytelling, terceiro
\end{textsample}
